% Live connection-distance distribution, snapshotted from the deployed
% network. n = 2274 connections across 396 active peers (snapshot
% 2026-05-25, source: the OTel-backed dashboard at nova.locut.us:3133).
%
% For each connection, ring distance = min(|loc_a - loc_b|, 1 - |loc_a - loc_b|).
% Linear histogram, 25 bins of width 0.02. The dashed curve is a
% Kleinberg 1/d reference with the same total mass.
%
% Coordinates are precomputed in cm (x: distance*22, y: count*5/450).
\begin{figure}[t]
\centering
\begin{tikzpicture}[
    bar/.style={fill=blue!22, draw=blue!55!black, line width=0.3pt},
    refcurve/.style={red!70!black, thick, dashed},
    axis/.style={->, >=Stealth, thick, black!75},
    tick/.style={black!75},
    label/.style={font=\scriptsize},
    gridline/.style={gray!20, very thin},
]
  % Background y-grid
    \draw[gridline] (0,1.1111) -- (11.0000,1.1111);
    \draw[gridline] (0,2.2222) -- (11.0000,2.2222);
    \draw[gridline] (0,3.3333) -- (11.0000,3.3333);
    \draw[gridline] (0,4.4444) -- (11.0000,4.4444);

  % Axes
  \draw[axis] (0,0) -- (11.4,0) node[right, label] {connection distance $d$};
  \draw[axis] (0,0) -- (0,5.4) node[above, label] {connections};

  % x-ticks
    \draw[tick] (0.0000,0) -- (0.0000,-0.1) node[below, label, yshift=-1pt] {$0.0$};
    \draw[tick] (2.2000,0) -- (2.2000,-0.1) node[below, label, yshift=-1pt] {$0.1$};
    \draw[tick] (4.4000,0) -- (4.4000,-0.1) node[below, label, yshift=-1pt] {$0.2$};
    \draw[tick] (6.6000,0) -- (6.6000,-0.1) node[below, label, yshift=-1pt] {$0.3$};
    \draw[tick] (8.8000,0) -- (8.8000,-0.1) node[below, label, yshift=-1pt] {$0.4$};
    \draw[tick] (11.0000,0) -- (11.0000,-0.1) node[below, label, yshift=-1pt] {$0.5$};

  % y-ticks
    \draw[tick] (0,0.0000) -- (-0.1,0.0000) node[left, label, xshift=-1pt] {0};
    \draw[tick] (0,1.1111) -- (-0.1,1.1111) node[left, label, xshift=-1pt] {100};
    \draw[tick] (0,2.2222) -- (-0.1,2.2222) node[left, label, xshift=-1pt] {200};
    \draw[tick] (0,3.3333) -- (-0.1,3.3333) node[left, label, xshift=-1pt] {300};
    \draw[tick] (0,4.4444) -- (-0.1,4.4444) node[left, label, xshift=-1pt] {400};

  % Bars (linear histogram, 25 bins of width 0.02)
    \draw[bar] (0.0000,0) rectangle (0.4400,4.5667);
    \draw[bar] (0.4400,0) rectangle (0.8800,3.2222);
    \draw[bar] (0.8800,0) rectangle (1.3200,2.7556);
    \draw[bar] (1.3200,0) rectangle (1.7600,1.9222);
    \draw[bar] (1.7600,0) rectangle (2.2000,1.5333);
    \draw[bar] (2.2000,0) rectangle (2.6400,1.1889);
    \draw[bar] (2.6400,0) rectangle (3.0800,1.0778);
    \draw[bar] (3.0800,0) rectangle (3.5200,1.0000);
    \draw[bar] (3.5200,0) rectangle (3.9600,0.9333);
    \draw[bar] (3.9600,0) rectangle (4.4000,0.9000);
    \draw[bar] (4.4000,0) rectangle (4.8400,0.6000);
    \draw[bar] (4.8400,0) rectangle (5.2800,0.6556);
    \draw[bar] (5.2800,0) rectangle (5.7200,0.5111);
    \draw[bar] (5.7200,0) rectangle (6.1600,0.4444);
    \draw[bar] (6.1600,0) rectangle (6.6000,0.3889);
    \draw[bar] (6.6000,0) rectangle (7.0400,0.3556);
    \draw[bar] (7.0400,0) rectangle (7.4800,0.4444);
    \draw[bar] (7.4800,0) rectangle (7.9200,0.4667);
    \draw[bar] (7.9200,0) rectangle (8.3600,0.3556);
    \draw[bar] (8.3600,0) rectangle (8.8000,0.5222);
    \draw[bar] (8.8000,0) rectangle (9.2400,0.3556);
    \draw[bar] (9.2400,0) rectangle (9.6800,0.4333);
    \draw[bar] (9.6800,0) rectangle (10.1200,0.2667);
    \draw[bar] (10.1200,0) rectangle (10.5600,0.1444);
    \draw[bar] (10.5600,0) rectangle (11.0000,0.2222);

  % 1/d reference (same total mass, bin midpoints, d_min clamp = 0.005).
  % Clip to the plot rectangle so the asymptote near d=0 doesn't escape.
  \begin{scope}
    \clip (0,0) rectangle (11.0,5.0);
    \draw[refcurve] plot coordinates {
      (0.2200,7.6056)
      (0.6600,3.8028)
      (1.1000,2.2244)
      (1.5400,1.5789)
      (1.9800,1.2244)
      (2.4200,1.0000)
      (2.8600,0.8456)
      (3.3000,0.7322)
      (3.7400,0.6467)
      (4.1800,0.5778)
      (4.6200,0.5233)
      (5.0600,0.4778)
      (5.5000,0.4389)
      (5.9400,0.4067)
      (6.3800,0.3789)
      (6.8200,0.3544)
      (7.2600,0.3322)
      (7.7000,0.3133)
      (8.1400,0.2967)
      (8.5800,0.2811)
      (9.0200,0.2678)
      (9.4600,0.2556)
      (9.9000,0.2433)
      (10.3400,0.2333)
      (10.7800,0.2244)
    };
  \end{scope}

  % Median guide line at d = 0.0819
  \draw[gray!50!black, thin, dashed] (1.8018,0) -- (1.8018,4.7)
       node[label, anchor=south west, gray!40!black] at (1.8018,4.7) {median $d = 0.082$};

  % Legend
  \node[label, anchor=north east, fill=white, draw=black!50,
        rounded corners=2pt, inner sep=4pt, align=left]
       at (10.9,4.9) {
       \textcolor{blue!55!black}{$\blacksquare$}\ empirical (live network)\\
       \textcolor{red!70!black}{- -}\ $1/d$ reference (Kleinberg)};

\end{tikzpicture}
\caption{Distribution of per-connection ring distances on the deployed
network. Each of the $n = 2274$ live connections across $N = 396$
active peers (snapshot taken 2026-05-25, source: the telemetry
collector at \texttt{nova.locut.us:3133}) contributes one data point;
bins are linear, width $0.02$. The dashed curve is a Kleinberg $1/d$
reference normalized to the same total mass. The empirical
distribution is monotonically decreasing in $d$ at short range and
stays above the $1/d$ reference in the long-range tail, consistent
with the active gap-targeting and Kleinberg-style acceptance of
\cref{sec:routing} steering each peer's neighborhood toward the
$1/d$ shape Kleinberg's analysis requires. Median connection
distance is $0.082$; the maximum half-ring distance is $0.5$.}
\label{fig:connection-distance}
\end{figure}
